\documentclass[12pt]{article}
\usepackage{graphicx,amsmath,amsfonts,amssymb,epsfig,euscript,enumerate}
\usepackage[T1]{fontenc}
\usepackage[utf8x]{inputenc}

\newtheorem{exercise}{Ejercicio}
\newcommand{\bej}{\begin{exercise}\rm}
\newcommand{\fej}{\end{exercise}}

\newcommand{\R}{\mathbb{R}}
\newcommand{\C}{\mathbb{C}}
\def\dt{\Delta t}
\def\dx{\Delta x}

\topmargin-2cm \vsize 29.5cm \hsize 21cm
\setlength{\textwidth}{16.75cm}\setlength{\textheight}{23.5cm}
\setlength{\oddsidemargin}{0.0cm}
\setlength{\evensidemargin}{0.0cm}

\begin{document}
\centerline{{\small Universidad de Buenos Aires - Facultad de Ciencias Exactas y Naturales - Depto. de Matemática}}
 
 \vskip 0.2cm
 \hrulefill
 \vskip 0.2cm

 \centerline{{\bf\Huge {\sc Elementos de Cálculo Numérico}}}
 \vskip 0.2cm
 \centerline{\ttfamily Primer Cuatrimestre 2026}
 \hrulefill

 \bigskip
 \centerline{\bf Parcial Modelo 1 (Guías 1 a 3)}
 \bigskip

%%% Ejercicio 1 — Guía 1: complejidad + aritmética de punto flotante %%%

\begin{exercise}[Evaluación eficiente de polinomios]
Sea $p(x) = a_0 + a_1 x + a_2 x^2 + \cdots + a_n x^n$.

\begin{enumerate}[(a)]
\item Considere el siguiente algoritmo que evalúa $p(x)$ calculando cada potencia $x^k$ desde cero:
\begin{verbatim}
S = a[0]
for k = 1 to n:
    xk = 1
    for j = 1 to k:
        xk = xk * x
    S = S + a[k] * xk
\end{verbatim}
Cuente el número total de multiplicaciones (operaciones \texttt{*}) y deduzca que la complejidad es $O(n^2)$.

\item El \textbf{método de Horner} reescribe el polinomio como
\[
p(x) = a_0 + x\bigl(a_1 + x\bigl(a_2 + \cdots + x(a_{n-1} + x\,a_n)\cdots\bigr)\bigr)
\]
y lo evalúa de adentro hacia afuera:
\begin{verbatim}
S = a[n]
for k = n-1 downto 0:
    S = S * x + a[k]
\end{verbatim}
Cuente el número total de multiplicaciones y sumas y deduzca la complejidad.

\item Sea $p(x) = x^3 - 3x^2 + 3x - 1$. Evalúe $p(1.01)$ paso a paso usando el método de Horner en aritmética de punto flotante con $t = 3$ dígitos decimales significativos (redondeo al más cercano).

\item Observe que $p(x) = (x-1)^3$. Evalúe $(1.01 - 1)^3$ con $t = 3$ dígitos y compare con el resultado del inciso anterior. Identifique en qué paso del método de Horner se produce pérdida catastrófica de dígitos significativos.
\end{enumerate}
\end{exercise}


%%% Ejercicio 2 — Guía 2: álgebra lineal numérica (SVD, cuadrados mínimos, leverage) %%%

\begin{exercise}[Cuadrados mínimos y puntos de alto leverage]
Sea $A \in \mathbb{R}^{m \times n}$ con $m > n$ y rango columna completo, y sea $A = U\Sigma V^T$ su SVD reducida ($U \in \mathbb{R}^{m \times n}$ con columnas ortonormales, $\Sigma \in \mathbb{R}^{n \times n}$ diagonal con $\sigma_1 \geq \cdots \geq \sigma_n > 0$, $V \in \mathbb{R}^{n \times n}$ ortogonal). La solución de cuadrados mínimos de $\min_x \|Ax - b\|_2$ es $\hat{x} = A^+ b$ y el vector de valores ajustados es $\hat{b} = A\hat{x}$.

\begin{enumerate}[(a)]
\item Demuestre que $\hat{x} = V\Sigma^{-1}U^T b$ y que $\hat{b} = Hb$, donde $H = UU^T \in \mathbb{R}^{m \times m}$ se llama \emph{hat matrix} (porque ``le pone el sombrero'' a $b$).

\item Demuestre que $H$ es un proyector ortogonal ($H^2 = H$, $H^T = H$). Deduzca que sus autovalores son $0$ o $1$, que las entradas diagonales satisfacen $0 \leq h_{ii} \leq 1$, y que $\sum_{i=1}^m h_{ii} = n$.

\textbf{Sugerencia para $h_{ii} \leq 1$:} use que $h_{ii} = e_i^T H e_i = \|He_i\|_2^2$ (¿por qué?) y que $\|H\|_2 \leq 1$.

\item Suponga que $b$ se perturba únicamente en la coordenada $i$: $\tilde{b} = b + \delta\, e_i$. Calcule el cambio en el valor ajustado $\hat{b}_i$ y concluya que $h_{ii}$ mide la fracción de la perturbación $\delta$ que se transmite al valor ajustado en el punto $i$.
En estadística, $h_{ii}$ se denomina el \emph{leverage} del punto $i$: si $h_{ii} \approx 1$, la observación $i$ determina casi completamente su propio valor ajustado, independientemente del resto de los datos.

\item Se ajusta una recta $\varphi(t) = c_0 + c_1 t$ a datos tomados en $t_1 = 0$, $t_2 = 1$, $t_3 = 2$, $t_4 = 9$, con la matriz de diseño
\[
A = \begin{pmatrix} 1 & 0 \\ 1 & 1 \\ 1 & 2 \\ 1 & 9 \end{pmatrix}.
\]
Calcule $(A^TA)^{-1}$ y los cuatro leverages $h_{ii} = a_i^T(A^TA)^{-1}a_i$, donde $a_i^T$ es la $i$-ésima fila de~$A$. Verifique que la suma de los leverages es $n = 2$. Determine qué punto tiene leverage más alto y qué porcentaje de una perturbación en ese dato se transmite directamente a su valor ajustado.
\end{enumerate}
\end{exercise}

%%% Ejercicio 3 — Guía 3: cuadratura + FFT %%%

\begin{exercise}[Beam gaussiano, regla de trapecios y FFT]
Considere la integral tipo convolución que modela un \emph{beam gaussiano}:
\[
I(x)=\int_{-L}^{L} f(\xi)\, \exp\!\left(-\frac{(x-\xi)^2}{\sigma^2}\right)\, d\xi.
\]
Sea una grilla uniforme $\xi_j=-L+jh$ con $h=\frac{2L}{N}$, $j=0,\ldots,N$, y sea $x_i=-L+ih$, $i=0,\ldots,N$.
\begin{enumerate}[(a)]
\item Discretice $I(x_i)$ con la regla de trapecios compuesta y escriba la suma resultante en términos de $f_j = f(\xi_j)$ y del kernel gaussiano evaluado en los nodos.
\item Defina $g_k = \exp\!\bigl(-{(kh)^2}/{\sigma^2}\bigr)$. Muestre que la suma del inciso anterior se puede escribir como
\[
I(x_i) \approx h \sum_{j=0}^{N} f_j \; g_{i-j},
\]
identificando el resultado como una convolución discreta (lineal) de las secuencias $\{f_j\}$ y $\{g_k\}$, salvo correcciones en los extremos del trapecio.

\item Para calcular los $N+1$ valores $I(x_0), \ldots, I(x_N)$ simultáneamente usando FFT:
\begin{enumerate}[i.]
\item Indique cómo extender las secuencias con ceros (\emph{zero-padding}) para que la convolución circular de la FFT reproduzca la convolución lineal deseada.
\item Describa los tres pasos del algoritmo (FFT directa, producto punto a punto, FFT inversa) y determine la complejidad total.
\end{enumerate}
\end{enumerate}
\end{exercise}

%%% Ejercicio 4 — Guía 3: interpolación %%%

\begin{exercise}[Interpolación inteligente: $\sqrt{x}\sin(x)$]
Se desea evaluar $f(x) = \sqrt{x}\,\sin(x)$ en puntos arbitrarios del intervalo $[0, \pi]$ usando interpolación polinomial en $n+1$ nodos equiespaciados $x_k = \frac{k\pi}{n}$, $k=0,\ldots,n$.

Se proponen dos estrategias:
\begin{itemize}
\item \textbf{Estrategia A:} Interpolar directamente $f(x) = \sqrt{x}\,\sin(x)$ con un polinomio $p_n(x)$ de grado $\leq n$.
\item \textbf{Estrategia B:} Interpolar solamente $g(x) = \sin(x)$ con un polinomio $q_n(x)$ de grado $\leq n$, y luego evaluar $\sqrt{x}\,q_n(x)$.
\end{itemize}

\begin{enumerate}[(a)]
\item El teorema del error de interpolación establece que si $\varphi \in C^{n+1}([a,b])$, entonces para cada $x$ existe $\xi$ tal que
\[
\varphi(x) - p_n(x) = \frac{\varphi^{(n+1)}(\xi)}{(n+1)!}\,\prod_{k=0}^{n}(x - x_k).
\]
Calcule $f'(x) = \frac{\sin x}{2\sqrt{x}} + \sqrt{x}\cos x$ y $f''(x)$. Observe que $f''(x)$ tiene un término proporcional a $x^{-3/2}$. Deduzca que $\|f^{(n+1)}\|_\infty$ en $[0, \pi]$ es no acotada (diverge cuando $x \to 0^+$). Concluya que el teorema del error no se puede aplicar directamente a la estrategia~A.

\item Calcule $g^{(n+1)}(x)$ para $g(x) = \sin(x)$ y demuestre que $\|g^{(n+1)}\|_\infty = 1$ para todo $n$.

\item Usando el teorema del error para la estrategia~B, demuestre que
\[
|f(x) - \sqrt{x}\,q_n(x)| \leq \frac{\sqrt{\pi}}{(n+1)!}\,\prod_{k=0}^{n}|x - x_k|
\]
para todo $x \in [0, \pi]$.

\item Para nodos equiespaciados con paso $h = \pi/n$, admita la cota estándar $\prod_{k=0}^{n}|x-x_k| \leq \frac{n!\, h^{n+1}}{4}$ (válida para $x$ entre dos nodos consecutivos). Deduzca que el error de la estrategia B satisface
\[
\|f - \sqrt{\cdot}\, q_n\|_\infty \leq \frac{\sqrt{\pi}\,\pi^{n+1}}{4\,(n+1)\,n^{n+1}}.
\]
Concluya que el error decrece más rápido que cualquier potencia de $1/n$.
\end{enumerate}
\end{exercise}

\end{document}
